\documentclass[10pt,a4paper]{amsart}
\usepackage[margin=19mm]{geometry}
\usepackage{amsmath,amssymb,mathtools}
\usepackage[scr=rsfs,cal=euler]{mathalfa}
\usepackage{xfrac}
\usepackage[unicode,hidelinks,bookmarks=false]{hyperref}
\newcommand{\ie}{i.e.\ }
\newcommand{\cf}{cf.\ }
\newcommand{\resp}{respectively\ }
\newcommand{\Subsection}{\subsection}
\providecommand{\nobreakdash}{\nobreak\hbox{-}}
\setlength{\emergencystretch}{2em}
\multlinegap0pt
\usepackage{amssymb}
\usepackage[shortlabels]{enumitem}
\newtheorem{proposition}{Proposition}
\newtheorem{lemma}{Lemma}
\newcommand{\Z}{\mathbb{Z}}
\title{Selfless reduced amalgamated free products and HNN extensions}
\author{David Gao, Srivatsav Kunnawalkam Elayavalli, Gregory Patchell, Lizzy Teryoshin}
\date{Source: arXiv:2604.06982v2 (CC BY 4.0)}
\begin{document}
\maketitle

\noindent\emph{Source and licence.} Selfless reduced amalgamated free products and HNN extensions. Version arXiv:2604.06982v2 (CC BY 4.0), distributed by arXiv under the Creative Commons Attribution 4.0 International licence, http://creativecommons.org/licenses/by/4.0/, as recorded in the arXiv licence metadata for this exact version and shown on the abstract page https://arxiv.org/abs/2604.06982v2. Full licence notice: \texttt{licenses/arxiv-2604.06982-CC-BY-4.0.txt}.\par\smallskip

\noindent\textit{Editorial scope.} Counted unit: the article's proposition hnn-univ-prop (source0.tex lines 220--230). The article's complete original proof of it is delivered with it (source0.tex lines 244--295, one \texttt{proof} environment of 52 lines, which the article places away from the statement). No mathematics is written, repaired or re-derived: the statement and the proof are the article's own lines, in the article's own notation, and no retained line is substituted or re-flowed.  Retained uncounted as the source-local context the proof needs: lines 234--242.  Nothing was omitted from any retained span, so the package carries no omission record.  No retained line contains a citation, so the wrapper carries no bibliography and no bibliography entry.  No retained line contains a figure environment, an image-inclusion command or drawing code, so no image is delivered and the package declares no asset dependency.  Editorial formatting only, in addition: an AMS \texttt{amsart} preamble that restates the article's own declarations and macros used by the retained lines, namely the environments \texttt{proposition}, \texttt{lemma} on separate counters, together with the standard packages those lines need.  The retained environments print in this document as Proposition 1, Lemma 1 in order of appearance, whereas the article prints the same items with the numbering of its own full document.  The complete native source of the article as distributed in the arXiv e-print for this exact version is bundled unchanged as \texttt{sources/198061/source0.tex}.
\begin{lemma}\label{lem: amal-free-state-zero}
    Consider the reduced amalgamated free product $P = P_1 \ast_R P_2$ with the canonical faithful state $\rho_P$. Let $x = y_0u_1y_1 \cdots y_{m-1}u_my_m$ with $m \ge 1$ and
    \begin{enumerate}
        \item $y_i \in P_1$ for all $0 \le i \le m$;
        \item $E_R(y_i) = 0$ for all $1 \le i \le m - 1$;
        \item $u_i \in P_2 \ominus R$.
    \end{enumerate}
    Then $\rho_P(x) = 0$.
\end{lemma}

\begin{proposition}\label{prop: hnn-univ-prop}
    Let $P = \mathrm{HNN}(A,\theta,B_1,B_{-1})$. Assume that $(Q,\rho_Q)$ is a $C^\ast$-algebra equipped with a faithful state, that $\pi: A \to Q$ is a state-preserving embedding, and that has a unitary $w_Q$. Suppose,
    \begin{enumerate}
        \item $\pi(\theta(x)) = w_Q\pi(x)w_Q^\ast$ for all $x\in B_1$;
        \item for all reduced $x = x_0w^{\varepsilon_1}x_1\cdots x_{n-1}w^{\varepsilon_n}x_n \in P,$ we have
        \begin{equation*}
            \rho_Q(\pi(x_0)w_Q^{\varepsilon_1}\pi(x_1)\cdots \pi(x_{n-1})w_Q^{\varepsilon_n}\pi(x_n)) = 0.
        \end{equation*}
    \end{enumerate}
    Then there exists a unique state-preserving *-homomorphism $\Tilde{\pi}:P\to Q$ extending $\pi$ and satisfying $\Tilde{\pi}(w) = w_Q.$
\end{proposition}

\begin{proof}[Proof of Proposition \ref{prop: hnn-univ-prop}]
    Since reduced words, together with scalars, span $P$, it suffices to show all reduced words have state zero in $P$. Clearly, this holds for elements of $A \ominus \mathbb{C}1$, so let $x = x_0w^{\varepsilon_1}x_1\cdots x_{n-1}w^{\varepsilon_n}x_n \in P$ be reduced. Now, recall that $P$ is a unital subalgebra of $C = ((A\otimes A)\rtimes \Z/2\Z) *_{(B_1\otimes B_{-1})}((B_1\otimes B_{-1})\rtimes \Z/2\Z)$. Recall that the copy of $A$ contained in $P$ is $A \otimes 1$ in $C$ and $w = uv$, where $u$ is the generator of the first copy of $\Z/2\Z$ and $v$ is the generator of the second copy of $\Z/2\Z$. Let $P_1 = (A\otimes A)\rtimes \Z/2\Z$, $P_2 = (B_1\otimes B_{-1})\rtimes \Z/2\Z$, and $R = B_1\otimes B_{-1}$. It then suffices to verify the conditions in Lemma \ref{lem: amal-free-state-zero}.

    For this, we proceed by induction on $n$ to prove that $x$, after combining terms, is indeed a product of the form given in Lemma \ref{lem: amal-free-state-zero}. Furthermore, the ending term $y_m$ equals $x_n$ if $\varepsilon_n = 1$ and equals $ux_n$ if $\varepsilon_n = -1$. Indeed, when $n = 1$, we have, if $\varepsilon_1 = 1$,
    \begin{equation*}
        x = x_0uvx_1 = (x_0u)(v)(x_1)
    \end{equation*}
    which is of the desired form. And, if $\varepsilon_1 = -1$,
    \begin{equation*}
        x = x_0vux_1 = (x_0)(v)(ux_1).
    \end{equation*}

    Now, assume the result holds for $n - 1$. There are 4 cases for the inductive step:
    \begin{enumerate}
        \item $\varepsilon_{n-1} = \varepsilon_n = 1$: By induction hypothesis,
        \begin{equation*}
            x_0w^{\varepsilon_1}x_1\cdots x_{n-2}w^{\varepsilon_{n-1}}x_{n-1} = y_0u_1y_1 \cdots y_{m-1}u_my_m
        \end{equation*}
        where the latter product satisfies the conditions in Lemma \ref{lem: amal-free-state-zero} and $y_m = x_{n-1}$. Thus, $y_mu = x_{n-1}u$ is orthogonal to $R$, so,
        \begin{equation*}
            x = y_0u_1y_1 \cdots y_{m-1}u_my_mwx_n = y_0u_1y_1 \cdots y_{m-1}u_m(y_mu)(v)(x_n)
        \end{equation*}
        is as required.
        \item $\varepsilon_{n-1} = 1$ and $\varepsilon_n = -1$: Again, by induction hypothesis,
        \begin{equation*}
            x_0w^{\varepsilon_1}x_1\cdots x_{n-2}w^{\varepsilon_{n-1}}x_{n-1} = y_0u_1y_1 \cdots y_{m-1}u_my_m
        \end{equation*}
        where the latter product satisfies the conditions in Lemma \ref{lem: amal-free-state-zero} and $y_m = x_{n-1}$. Furthermore, by definition of reduced words, $x_{n-1} \perp B_1$, so $y_m = x_{n-1} \otimes 1$ is orthogonal to $R = B_1 \otimes B_{-1}$. Thus,
        \begin{equation*}
            x = y_0u_1y_1 \cdots y_{m-1}u_my_mw^{-1}x_n = y_0u_1y_1 \cdots y_{m-1}u_m(y_m)(v)(ux_n)
        \end{equation*}
        is as required.
        \item $\varepsilon_{n-1} = -1$ and $\varepsilon_n = 1$: By induction hypothesis,
        \begin{equation*}
            x_0w^{\varepsilon_1}x_1\cdots x_{n-2}w^{\varepsilon_{n-1}}x_{n-1} = y_0u_1y_1 \cdots y_{m-1}u_my_m
        \end{equation*}
        where the latter product satisfies the conditions in Lemma \ref{lem: amal-free-state-zero} and $y_m = ux_{n-1}$. Furthermore, by definition of reduced words, $x_{n-1} \perp B_{-1}$, so $y_mu = ux_{n-1}u = 1 \otimes x_{n-1}$ is orthogonal to $R = B_1 \otimes B_{-1}$. Thus,
        \begin{equation*}
            x = y_0u_1y_1 \cdots y_{m-1}u_my_mwx_n = y_0u_1y_1 \cdots y_{m-1}u_m(y_mu)(v)(x_n)
        \end{equation*}
        is as required.
        \item $\varepsilon_{n-1} = \varepsilon_n = -1$: By induction hypothesis,
        \begin{equation*}
            x_0w^{\varepsilon_1}x_1\cdots x_{n-2}w^{\varepsilon_{n-1}}x_{n-1} = y_0u_1y_1 \cdots y_{m-1}u_my_m
        \end{equation*}
        where the latter product satisfies the conditions in Lemma \ref{lem: amal-free-state-zero} and $y_m = ux_{n-1}$. Thus, $y_m$ is orthogonal to $R$, so,
        \begin{equation*}
            x = y_0u_1y_1 \cdots y_{m-1}u_my_mw^{-1}x_n = y_0u_1y_1 \cdots y_{m-1}u_m(y_m)(v)(ux_n)
        \end{equation*}
        is as required.
    \end{enumerate}
\end{proof}

\end{document}
